\section{Direct quartic strict vanishing in the first weight residue}
\label{sec:cubic-quartic-strict-vanishing}

The degree-twelve class certificate AC-G129 has a direct consequence
over the original field \(K\).  An exact unit calculation eliminates
both nontrivial cubic characters for each of the two primitive
quartic characters modulo five.  The passage to the original
discrete coefficients is an explicit uniformizer argument; no
automorphic seed, auxiliary host, or deformation ring is a premise.
Only the residue \(267\) modulo \(268\) is treated here.

\paragraph{Characters, coefficients, and the precise strict condition.}
Put
\[
 p=269,\qquad K=\mathbb Q(\sqrt{-3}),\qquad t^3=p,\qquad
 z=\zeta_5,\qquad M=K(t,z),\qquad L=\mathbb Q(t,z).
 \tag{QSV.1}
\]
Fix either generator
\(\chi:G_K\longrightarrow\mu_3\) of \(K(t)/K\).
Let \(\psi\) be the primitive Dirichlet character modulo five with
\(\psi(2)=i\), where the \(269\)-adic fourth root \(i\) reduces to
\(82\).  Its arithmetic Galois avatar is \(\Psi\): if
\(g(z)=z^a\), then \(\Psi(g)=\psi(a)\).
Write \(c_p\) for the actual action on \(p\)-power roots of unity.
Thus \(c_p\) takes arithmetic Frobenius at an unramified prime \(q\)
to \(q\); no inverse Tate convention is used.

For \(\sigma\in\{1,-1\}\), \(j\in\{1,2\}\), and
\(e\ge1\) with \(e\equiv267\pmod{268}\), put
\[
             a_{j,\sigma,e}=\chi^{-j}\Psi^{-\sigma}c_p^{-e}.
 \tag{QSV.2}
\]
Let \(k\) be any finite field of characteristic \(p\) containing
the reduced character values.  For either a finite discrete
coefficient module or the discrete module \(E/O\) below, use exactly
\[
 \begin{split}
 \operatorname{Sel}_{\mathrm{str}}(K,V)=\ker\biggl(
 H^1_{\mathrm{cts}}(G_K,V)&\longrightarrow
       \prod_{v\mid p}H^1_{\mathrm{cts}}(G_{K_v},V)\\
       &\hspace{8mm}\times
       \prod_{v\nmid p,\infty}H^1_{\mathrm{cts}}(I_v,V)
                         \biggr).
 \end{split}
 \tag{QSV.3}
\]
The condition at \(p\) is zero in the full local cohomology group.
Away from \(p\) it is unramifiedness.  Complex places contribute
nothing.  In particular no Bloch--Kato finite condition replaces the
local zero condition in (QSV.3).

\begin{theorem}[AC-G130-residual: both quartic and cubic residual sectors]
\label{thm:ACG130residual}
For all data in (QSV.1)--(QSV.3),
\[
       \operatorname{Sel}_{\mathrm{str}}
                (K,k(\bar a_{j,\sigma,e}))=0
             \qquad(j=1,2,\ \sigma=1,-1).
 \tag{QSV.4}
\]
The class input is precisely AC-G129: the nontrivial cubic part of
\(\operatorname{Cl}(M)[269]\) is zero, as deduced there from the
certified prime-to-\(269\) class number of \(L\).
The unit input is the single verified global unit \(u_4\) from that
certificate, with the exact local computation below.
No class assertion about another quartic reflection field is used.
\end{theorem}

\begin{proof}
\proofstep{1. Exact inverse-character Kummer descent.}
Reduction of \(e\equiv-1\pmod{p-1}\) gives
\[
 \bar a_{j,\sigma,e}=\bar\chi^{-j}\bar\Psi^{-\sigma}\bar c_p,
 \qquad
 \eta_{j,\sigma}:=\bar a_{j,\sigma,e}/\bar c_p
                  =\bar\chi^{-j}\bar\Psi^{-\sigma}.
 \tag{QSV.5}
\]
The cubic and quartic extensions of \(K\) have coprime degrees and
are nontrivial.  Their product character is faithful on
\[
       \Delta=\operatorname{Gal}(M/K)=C_3\times C_4.
 \tag{QSV.6}
\]
Thus \(M\), rather than a field with an additional \(\mu_p\), is
the splitting field of \(\eta_{j,\sigma}\).
The two factors in the normal closure over \(\mathbb Q\) are
\(S_3\) and \(C_4\): their possible quadratic intersection would
be \(K\), which is ramified at three whereas \(\mathbb Q(z)\) is not.

Let \(S\) be the primes of \(M\) above \(p\), and set
\[
 \begin{split}
 V_{M,S}&=\{x\in M^\times:
       p\mid\operatorname{ord}_w(x)\text{ for }w\nmid p\}/M^{\times p},\\
 \mathcal L_M&=\bigoplus_{w\mid p}M_w^\times/M_w^{\times p}.
 \end{split}
 \tag{QSV.7}
\]
Prime-to-\(p\) inflation--restriction and the Kummer sequence give
\[
 \operatorname{Sel}_{\mathrm{str}}(K,k(\bar a_{j,\sigma,e}))
 \simeq
 e_{\eta_{j,\sigma}^{-1}}\ker\left(
       V_{M,S}\otimes_{\mathbb F_p}k
                \longrightarrow\mathcal L_M\otimes_{\mathbb F_p}k
                              \right).
 \tag{QSV.8}
\]
Here the natural \(\Delta\)-action on
\(H^1(M,\mu_p)=M^\times/M^{\times p}\) is twisted by
\(\eta_{j,\sigma}\).  Invariance therefore selects the
\emph{inverse} character
\(\eta_{j,\sigma}^{-1}=\bar\chi^j\bar\Psi^\sigma\).
The same inflation--restriction argument applies to every local
decomposition group and every away inertia group, since its
quotient has order dividing twelve.  It preserves the actual
zero restrictions.  Away from \(p\), Kummer unramifiedness means
valuation divisible by \(p\); at \(p\), zero means a local
\(p\)th power.  Every \(p\)-adic completion in (QSV.7) is retained.

\proofstep{2. Strictness puts the class in the ordinary unit sector.}
Define the all-primes valuation kernel
\[
 V_M=\{x\in M^\times:
          p\mid\operatorname{ord}_w(x)
              \text{ for every finite }w\}/M^{\times p}.
 \tag{QSV.9}
\]
If the localization in (QSV.8) is zero, its valuations at the
primes above \(p\) are zero modulo \(p\) as well.  Such a class
lies in \(V_M\otimes k\).  This assertion also holds for arbitrary
linear combinations after extension of scalars: \(V_M\) is the
kernel of the remaining valuation map on \(V_{M,S}\), and field
extension is exact.

Put \(E_M=\mathcal O_M^\times\).  The elementary ideal-factorization
sequence is
\[
       0\longrightarrow E_M/E_M^p
        \longrightarrow V_M
        \longrightarrow\operatorname{Cl}(M)[p]
        \longrightarrow0.
 \tag{QSV.10}
\]
Indeed, for \(x\in V_M\) write \((x)=\mathfrak a^p\);
the final map is \([x]\mapsto[\mathfrak a]\).  Its kernel is a
global unit modulo a \(p\)th power, and every \(p\)-torsion ideal
class has such a representative.  All maps are \(\Delta\)-equivariant.
The \(\eta_{j,\sigma}^{-1}\)-part of the last term is zero by
AC-G129, because its cubic character is nontrivial.
Consequently every class in (QSV.8) comes from the corresponding
ordinary global-unit sector.  No \(S\)-class correction or
unproved \(S\)-unit rank equality is required.

\proofstep{3. The exact two-dimensional unit modules.}
There is no \(\mu_p\) in \(M\), since \([M:\mathbb Q]=24\) whereas
\([\mathbb Q(\mu_p):\mathbb Q]=268\).
The twelve complex places of \(M\) form a free transitive
\(\Delta\)-set: choose the embedding over one fixed complex
embedding of the imaginary quadratic field \(K\).
The equivariant logarithmic unit theorem identifies the rational
unit character with the regular character of \(\Delta\) minus
the trivial character.

This character calculation gives the same multiplicities after
reduction modulo \(p\).  To see the integral point, extend
\(\mathbb Z_p\) unramifiedly to contain the twelfth roots of unity.
The character idempotents have denominators dividing twelve and
therefore preserve the unit lattice.  Their ranks can be read in
characteristic zero and are unchanged after reduction.  The
absence of \(p\)-power torsion identifies that reduction with
\(E_M/E_M^p\).  This is a multiplicity argument, not an assumption
that a computed list of units is fundamental.

Let \(b\) generate \(C_4\) by \(b(z)=z^2\), fixing \(K\) and \(t\);
write \(\tau\) for a generator of \(C_3\).  The characters
\(\bar\Psi^{\pm1}\) are already defined over \(\mathbb F_p\), since
\(82^2=-1\pmod p\).  In the following character algebra, \(i\)
denotes its residue \(82\in\mathbb F_p\).  Define
\[
 \begin{split}
 e_\sigma&=\frac14\sum_{r=0}^3 i^{-\sigma r}b^r,\\
 e_{\mathrm{nt}}&=1-\frac13(1+\tau+\tau^2),\\
 W_\sigma&=e_{\mathrm{nt}}e_\sigma(E_M/E_M^p),
                  \qquad \sigma=1,-1 .
 \end{split}
 \tag{QSV.11}
\]
Each \(W_\sigma\) has dimension two over \(\mathbb F_p\).
The polynomial \(X^2+X+1\) is irreducible over \(\mathbb F_p\),
because \(p\equiv2\pmod3\).  Thus \(W_\sigma\) is the simple
nontrivial two-dimensional \(\mathbb F_p[C_3]\)-module.
After extension to \(k\), its two one-dimensional constituents
are exactly the characters
\(\bar\chi^j\bar\Psi^\sigma\), \(j=1,2\).

\proofstep{4. A global unit and both exact local projectors.}
Use the global unit \(u=u_4\in L^\times\) certified in
(QCC.3) of AC-G129.  To identify it without a choice of numerical
roots, write \(u=(P_0(z)+tP_1(z)+t^2P_2(z))/3\), where the
ascending coefficient vectors are
\[
 \begin{split}
 P_0:\;&(123029571,373382222,405083337,174314690),\\
 P_1:\;&(-39479967,-52428731,-20952996,11450068),\\
 P_2:\;&(3164085,-837481,-6474546,-5956867).
 \end{split}
 \tag{QSV.12}
\]
Its exact monic integral degree-twelve relation has constant one;
the verified rational polynomial annihilation proves that \(u\)
and \(u^{-1}\) are integral.  This is the previously verified
single unit, not a new assertion about a complete unit basis.

At \(p\), \(K_p/\mathbb Q_p\) is the unramified quadratic
extension.  It contains all fifth roots of unity because
\(5\mid p^2-1\), and \(t^3=p\) is Eisenstein over \(K_p\).
Thus every completion of \(M\) above \(p\) has absolute
ramification degree three and residue degree two; \(t\) is a
uniformizer.  Choose one completion and identify its residue
field with
\[
       \mathbb F_{p^2}=\mathbb F_p[s]/(s^2+3),\qquad
       z\bmod t=233+14s .
 \tag{QSV.13}
\]
The polynomial \(s^2+3\) is irreducible, and the displayed element
is a primitive fifth root.  These identities are checked by exact
modular arithmetic in the certificate.
Reducing (QSV.12), the expansion of \(b^r(u)\) is
\(A_r+B_rt+O(t^2)\), with the following pairs; a pair \((a,b)\)
denotes \(a+bs\).
\[
 \begin{array}{c|c|c|c}
 r&A_r&B_r&B_r/A_r\\ \hline
 0&(170,68)&(259,4)&(192,15)\\
 1&(233,255)&(257,207)&(77,181)\\
 2&(170,201)&(259,265)&(192,254)\\
 3&(233,14)&(257,62)&(77,88)
 \end{array}
 \tag{QSV.14}
\]
Every \(A_r\) is nonzero.  The integers representing the two
projectors in (QSV.11) are
\[
      (w_{+,r})=(202,114,67,155),\qquad
      (w_{-,r})=(202,155,67,114).
 \tag{QSV.15}
\]
Form actual global units
\[
 a_\sigma=\prod_{r=0}^3 b^r(u)^{w_{\sigma,r}},\qquad
 v_\sigma=\frac{a_\sigma^3}
                 {a_\sigma\,\tau(a_\sigma)\,\tau^2(a_\sigma)} .
 \tag{QSV.16}
\]
Their classes satisfy
\([v_\sigma]=3e_{\mathrm{nt}}e_\sigma[u]\in W_\sigma\).
All exponents in (QSV.16) are integers; division by four and
three occurs only in the prime-to-\(p\) character algebra.

The cubic norm has zero linear term, since
\(1+\omega+\omega^2=0\) and \(\tau(t)=\omega t\).
Therefore the exact local expansions of (QSV.16) are
\[
 \begin{split}
 v_+&=1+
       3\sum_{r=0}^3w_{+,r}(B_r/A_r)t+O(t^2)
          =1+221st+O(t^2),\\
 v_-&=1+
       3\sum_{r=0}^3w_{-,r}(B_r/A_r)t+O(t^2)
          =1+93st+O(t^2).
 \end{split}
 \tag{QSV.17}
\]
The verifier also multiplies these expressions directly in the
dual-number ring over \(\mathbb F_{p^2}\), using
\(\omega=(-1+s)/2\), and checks the norm correction before
reading off the coefficients.  Both computations are exact.

Neither local unit in (QSV.17) is a \(p\)th power.  Indeed, if
\(v=y^p\) reduces to one, then \(y\) also reduces to one because
the residue multiplicative group has order \(p^2-1\), prime to
\(p\).  Write \(y=1+x\), with \(v_t(x)\ge1\).
For \(1\le a<p\), the binomial coefficient \(\binom pa\) has
\(t\)-valuation three.  Hence every term of
\((1+x)^p-1\) has valuation at least
\(\min(3+v_t(x),p\,v_t(x))\ge4\).
It cannot have either nonzero linear coefficient in (QSV.17).

\proofstep{5. Injectivity and the residual conclusion.}
Localization of global units into the sum over all \(p\)-adic
completions is \(\Delta\)-equivariant.  Equation (QSV.17) makes
its restriction to \(W_\sigma\) nonzero.  Simplicity of
the two-dimensional \(\mathbb F_p[C_3]\)-module \(W_\sigma\)
therefore makes this restriction injective.
Extension of scalars to \(k\) preserves injectivity, including
each of its two cubic-character lines.
By Steps1--2, every residual strict class belongs to one of those
lines and has zero localization.  It is consequently zero.
This proves (QSV.4) for both \(j\) and both \(\sigma\).
\end{proof}

\begin{theorem}[AC-G130-strict: original coefficients and prime-character lifts]
\label{thm:ACG130strict}
Let \(E/\mathbb Q_p\) be any finite extension with ring \(O\),
uniformizer \(\varpi\), and residue field \(k\), containing the
values of \(\chi\).  Keep this original ring \(O\).
Let \(\omega_p\) be the Teichmuller lift of \(\bar c_p\), let
\(r\in\mathbb Z\), and let
\(\kappa:G_K\to1+\mathfrak m_O\) be any continuous character
unramified at every finite place away from \(p\).
For \(e\ge1\) with \(e+r\equiv267\pmod{268}\), put
\[
 \begin{split}
 \psi_{\sigma,r,\kappa}&=\Psi^\sigma\omega_p^r\kappa,\\
 a_{j,\sigma,r,\kappa,e}
    &=\chi^{-j}\psi_{\sigma,r,\kappa}^{-1}c_p^{-e},\\
 D_{j,\sigma,r,\kappa,e}
    &=(E/O)(a_{j,\sigma,r,\kappa,e}).
 \end{split}
 \tag{QSV.18}
\]
Then
\[
       \operatorname{Sel}_{\mathrm{str}}
                (K,D_{j,\sigma,r,\kappa,e})=0
                   \qquad(j=1,2,\ \sigma=1,-1).
 \tag{QSV.19}
\]
In particular both its \(O\)-length and its normalized
\(O\)-length divided by \(e_O=\operatorname{length}_O(O/pO)\)
are zero.  No hypothesis \(e_O=1\), finite order of \(\kappa\),
or extension of the original coefficient ring is required.
\end{theorem}

\begin{proof}
\proofstep{1. The residual group and all local invariants.}
Write \(D=D_{j,\sigma,r,\kappa,e}\) and \(a\) for its character.
Its reduction is
\(\bar\chi^{-j}\bar\Psi^{-\sigma}\bar c_p^{-(e+r)}\).
The preceding theorem makes its strict residual group zero.
Multiplication by the original uniformizer gives an exact
sequence of discrete continuous modules
\[
   0\longrightarrow k(\bar a)
    \longrightarrow D\xrightarrow{\varpi}D\longrightarrow0,
   \qquad
   \bar u\longmapsto u/\varpi+O .
 \tag{QSV.20}
\]
This is a uniformizer kernel, not \(p\)-torsion over a possibly
ramified coefficient ring.

At \(p\), the cubic character has tame inertia image of order
three, whereas \(\bar c_p\) has inertia order dividing \(268\);
\(\Psi\) is unramified and \(\bar\kappa=1\).
Raising an inertia element with nontrivial cubic value to the
\(268\)th power kills its rational cyclotomic value and preserves
a nontrivial cubic value.  Thus some \(g\) in this inertia has
\(\bar a(g)\ne1\), and \(a(g)-1\) is a unit of \(O\).
It follows that
\[
                  D^{G_{K_p}}=0,\qquad D^{G_K}=0 .
 \tag{QSV.21}
\]
At a finite place \(v\nmid p\), the characters \(c_p,\omega_p\)
and \(\kappa\) are unramified.  The remaining inertia character
\(\chi^{-j}\Psi^{-\sigma}\) has finite order dividing twelve,
prime to \(p\).  If it is trivial, \(D^{I_v}=D\) is
\(\varpi\)-divisible.  Otherwise a nontrivial prime-to-\(p\)
root of unity differs from one by a unit of \(O\), so \(D^{I_v}=0\).
Consequently every local invariant module in (QSV.3) is
\(\varpi\)-divisible.  This checks the places above three and
five as well as every other away place.

\proofstep{2. The exact original-uniformizer lift.}
The global cohomology sequence of (QSV.20), using (QSV.21), gives
\[
       H^1(G_K,k(\bar a))\simeq H^1(G_K,D)[\varpi].
 \tag{QSV.22}
\]
For each local group \(J\) in (QSV.3), divisibility of \(D^J\)
makes \(H^1(J,k(\bar a))\to H^1(J,D)\) injective.
Thus the lift in (QSV.22) of a strict uniformizer-torsion class
is a residual strict class, and must be zero.

Equivalently, if \(\varpi z=dm\), choose \(n\in D\) with
\(\varpi n=m\); then \(z-dn\) takes values in \(D[\varpi]\).
If this residual cocycle is locally \(du\) in \(D\), then
\(\varpi u\in D^J\).  Choose \(u_0\in D^J\) with
\(\varpi u_0=\varpi u\).  Now \(u-u_0\in D[\varpi]\) and
\(du=d(u-u_0)\), proving local injectivity and preserving the
literal zero condition.  All maps and coboundaries are continuous.

\proofstep{3. No nonzero class can survive.}
Every continuous cocycle from the compact profinite group \(G_K\)
to the discrete module \(E/O\) has finite image.  One power of
\(\varpi\) kills every value of that cocycle.  Hence every
cohomology class is uniformizer-primary torsion.
A nonzero strict class would have a least annihilating power,
and its penultimate multiple would be nonzero strict
\(\varpi\)-torsion.  Step2 excludes this.  Equation (QSV.19)
follows without a finiteness or Iwasawa hypothesis.
\end{proof}

\paragraph{A finite Dirichlet family, with exact counting scope.}
Fix \(e\ge1\) and \(m\ge1\).  Among Dirichlet characters of
conductor dividing \(p^m\), the tame character is specified by
one exponent modulo \(268\), and the wild character has
\(p^{m-1}\) possibilities.  Taking the unique
\(r\equiv267-e\pmod{268}\), each of the two primitive quartic
labels at five therefore gives \(p^{m-1}\) characters
\(\psi^\sigma\omega_p^r\kappa\).  Altogether there are
\[
                           2\cdot269^{m-1}
 \tag{QSV.23}
\]
distinct sufficient Dirichlet labels of conductor dividing
\(5\cdot269^m\).  They remain distinct on \(G_K\):
a quotient trivial there would be \(1\) or \(\eta_K\);
the latter is ramified at three whereas these quotients are not.
The two quartic choices differ at five and cannot be identified
by a \(p\)-power conductor character.
Their parity is
\(\psi^\sigma(-1)(-1)^r=(-1)^{r+1}=(-1)^e\), because a
wild odd-prime-order character is trivial on \(-1\).
Every coefficient realization containing the chosen values is
covered by (QSV.19), with its own unchanged \(O\).
This is a count of sufficient character labels, not a family-average
or original arithmetic counting-budget theorem.

\paragraph{Certificate and source scope.}
The class input is the fully reviewed AC-G129 analytic
interval, exact finite-unit and saturation certificate; no
provisional PARI class output is substituted for it.
The local script
\texttt{verify\_quartic\_local\_projectors.py} reads the
verified unit numerator array, works in
\(\mathbb F_{269}[s]/(s^2+3)\), and checks both products in
the dual-number ring.  Its output is
\texttt{quartic\_local\_projectors\_certificate.json}.
It does not assume a complete unit system or a full class-group
structure.  The Kummer sequence is the same retained Stacks
Project tag03PK and unit/Picard interface used in AC-G103 and
AC-G114.  The exact inverse-character descent, unit multiplicity,
and original-coefficient arguments have been rederived above.

Geometrically, (QSV.8) is the indicated character part of Kummer
torsors over \(\operatorname{Spec}\mathcal O_M[1/p]\), trivial
over every \(p\)-adic completion.  Their valuation classes extend
to the all-primes kernel, the class-sector certificate removes
its ideal obstruction, and the unit localization is injective.
The argument concerns these actual torsor classes and their
explicit uniformizer lifts.  It asserts no new automorphic
support, mixed Bloch--Kato vanishing, \(e=133\) certificate,
other-prime conclusion, or H04.4 window inequality.
